%==============================================================================
%  CAHIER DES CHARGES
%  Evolution de l'architecture reseau : du routage par port
%  vers l'identification par adresse MAC
%  Stage de 45 jours (7 semaines)
%  Compilation : latexmk -pdf Cahier_des_Charges_MAC_Forwarding.tex
%==============================================================================
\documentclass[11pt,a4paper]{article}

%------------------------- Encodage et langue ---------------------------------
\usepackage[utf8]{inputenc}
\usepackage[T1]{fontenc}
\usepackage[french]{babel}
\usepackage{newtxtext,newtxmath}
\usepackage{microtype}

%------------------------- Mise en page ---------------------------------------
\usepackage[a4paper,top=2.4cm,bottom=2.2cm,left=2.2cm,right=2.2cm,headheight=15pt]{geometry}
\usepackage{fancyhdr}
\usepackage{titlesec}
\usepackage{lastpage}
\usepackage{setspace}
\setstretch{1.06}

%------------------------- Tableaux et listes ---------------------------------
\usepackage{array}
\usepackage{tabularx}
\usepackage{booktabs}
\usepackage{longtable}
\usepackage{colortbl}
\usepackage{ragged2e}
\usepackage{enumitem}
\usepackage{float}
\usepackage{caption}

%------------------------- Graphiques -----------------------------------------
\usepackage{xcolor}
\usepackage{tikz}
\usetikzlibrary{arrows.meta,positioning,calc}
\usepackage{pgfgantt}
\usepackage[most]{tcolorbox}
\usepackage{listings}

%------------------------- Couleurs de la charte ------------------------------
\definecolor{primary}{HTML}{00205B}
\definecolor{accent}{HTML}{0F7FA5}
\definecolor{soft}{HTML}{E8EEF5}
\definecolor{warn}{HTML}{B3541E}
\definecolor{okgreen}{HTML}{1F7A4D}
\definecolor{gris}{HTML}{5A5A5A}

%------------------------- Hyperliens -----------------------------------------
\usepackage[hidelinks]{hyperref}
\hypersetup{
  colorlinks=true, linkcolor=primary, urlcolor=accent,
  pdftitle={Cahier des charges - Evolution de l'architecture reseau vers l'identification par adresse MAC},
  pdfauthor={Airbus - Stage Infrastructure et Securite Reseau},
  pdfsubject={Cahier des charges technique},
  pdfkeywords={VLAN, MAC, 802.1X, MAB, RADIUS, EtherChannel, Palo Alto, NetApp, Centreon}
}

%------------------------- Styles de titres -----------------------------------
\titleformat{\section}{\normalfont\large\sffamily\bfseries\color{primary}}{\thesection.}{0.7em}{}
\titlespacing*{\section}{0pt}{14pt}{7pt}
\titleformat{\subsection}{\normalfont\normalsize\sffamily\bfseries\color{accent!85!black}}{\thesubsection}{0.7em}{}
\titlespacing*{\subsection}{0pt}{10pt}{4pt}

%------------------------- En-tetes / pieds de page ---------------------------
\pagestyle{fancy}
\fancyhf{}
\renewcommand{\headrulewidth}{0.4pt}
\renewcommand{\footrulewidth}{0.4pt}
\fancyhead[L]{\footnotesize\textcolor{gris}{Cahier des charges --- Identification par adresse MAC}}
\fancyhead[R]{\footnotesize\textcolor{gris}{v1.0}}
\fancyfoot[L]{\footnotesize\textcolor{gris}{Airbus --- Diffusion interne}}
\fancyfoot[R]{\footnotesize\textcolor{gris}{Page \thepage\ / \pageref{LastPage}}}

%------------------------- Colonnes personnalisees ----------------------------
\newcolumntype{Y}{>{\RaggedRight\arraybackslash}X}
\newcolumntype{L}[1]{>{\RaggedRight\arraybackslash}p{#1}}
\newcolumntype{C}[1]{>{\Centering\arraybackslash}p{#1}}
\newcolumntype{K}{>{\bfseries\ttfamily\footnotesize}l}

\captionsetup{font=small,labelfont={bf,color=primary},skip=5pt}
\renewcommand{\arraystretch}{1.18}

%------------------------- Boites ---------------------------------------------
\newtcolorbox{notebox}[1][Note technique]{
  enhanced, breakable, colback=soft, colframe=primary,
  boxrule=0.7pt, arc=2pt, left=7pt, right=7pt, top=5pt, bottom=5pt,
  fonttitle=\bfseries\sffamily\small, title={#1}, coltitle=white}

\newtcolorbox{alertbox}[1][Point de vigilance]{
  enhanced, breakable, colback=warn!6, colframe=warn,
  boxrule=0.7pt, arc=2pt, left=7pt, right=7pt, top=5pt, bottom=5pt,
  fonttitle=\bfseries\sffamily\small, title={#1}, coltitle=white}

%------------------------- Style des listings ---------------------------------
\lstdefinestyle{cli}{
  basicstyle=\ttfamily\scriptsize,
  keywordstyle=\color{primary}\bfseries,
  commentstyle=\color{okgreen}\itshape,
  numbers=left, numberstyle=\tiny\color{gris}, numbersep=6pt,
  frame=single, rulecolor=\color{black!25},
  backgroundcolor=\color{black!3},
  breaklines=true, showstringspaces=false,
  columns=fullflexible, keepspaces=true,
  xleftmargin=12pt, framexleftmargin=12pt, aboveskip=6pt, belowskip=4pt,
  morecomment=[l]{!}, morecomment=[l]{\#},
  morekeywords={interface,switchport,authentication,dot1x,mab,radius,server,ip,dhcp,
  snooping,arp,inspection,spanning-tree,port-channel,channel-group,channel-protocol,
  lacp,vlan,igroup,create,lun,map,show,aaa,address,key,trunk,native,allowed,mode,
  access,event,fail,action,authorize,violation,maximum,portfast,bpduguard,enable,
  load-balance,verify,source,nonegotiate}
}
\lstset{style=cli}

%------------------------- Macros ---------------------------------------------
\newcommand{\ttb}[1]{\texttt{\small #1}}
\newcommand{\rid}[1]{\textcolor{primary}{\textbf{\small #1}}}

\begin{document}

%==============================================================================
% PAGE DE GARDE
%==============================================================================
\begin{titlepage}
\begin{center}
\begin{tikzpicture}[remember picture,overlay]
  \fill[primary] ($(current page.north west)+(0,-3.0cm)$) rectangle ($(current page.north east)+(0,-3.4cm)$);
  \fill[accent]  ($(current page.south west)+(0,2.6cm)$) rectangle ($(current page.south east)+(0,2.68cm)$);
\end{tikzpicture}

\vspace*{0.5cm}
{\sffamily\Large\bfseries\color{primary} AIRBUS}\\[2pt]
{\sffamily\small\color{gris} Direction des Systèmes d'Information --- Infrastructure et Sécurité Réseau}

\vspace{2.4cm}

{\sffamily\Large\bfseries\color{gris} CAHIER DES CHARGES}

\vspace{1.2cm}

{\sffamily\huge\bfseries\color{primary}
Évolution de l'architecture réseau :\\[6pt] du routage par port vers\\[6pt] l'identification par adresse MAC}

\vspace{0.8cm}
{\large\itshape\color{gris} Segmentation VLAN, répartition de charge par hachage MAC\\
et sécurisation de l'infrastructure --- validation sur maquette}

\vspace{2.0cm}

\begin{tabular}{@{}>{\bfseries\sffamily\color{primary}}l@{\hspace{1.2em}}l@{}}
\toprule
Référence & CDC-RES-MAC-2026-001 \\
Version & 1.0 --- Pour validation \\
Classification & Diffusion interne \\
Date d'édition & \today \\
Durée du stage & 45 jours calendaires (7 semaines) \\
\midrule
Auteur & \makebox[6cm][l]{\dotfill} \\
Tuteur entreprise & \makebox[6cm][l]{\dotfill} \\
Tuteur académique & \makebox[6cm][l]{\dotfill} \\
Période & \makebox[6cm][l]{\dotfill} \\
\bottomrule
\end{tabular}

\vfill
{\footnotesize\color{gris} Document de travail --- Ne pas diffuser à l'extérieur de l'organisation.}
\end{center}
\end{titlepage}

%==============================================================================
% FICHE DE SUIVI ET SOMMAIRE
%==============================================================================
\section*{Fiche de suivi du document}
\addcontentsline{toc}{section}{Fiche de suivi du document}

\begin{tabular}{C{1.4cm} C{2.4cm} C{2.8cm} L{6.7cm}}
\toprule
\rowcolor{soft}\textbf{Version} & \textbf{Date} & \textbf{Rédacteur} & \textbf{Nature de la modification} \\
\midrule
0.5 & \dotfill & \dotfill & Trame initiale et étude de l'existant. \\
1.0 & \today & \dotfill & Version complète soumise à validation. \\
\bottomrule
\end{tabular}

\vspace{10pt}

\begin{tabular}{L{3.1cm} L{4.9cm} L{2.7cm} L{2.6cm}}
\toprule
\rowcolor{soft}\textbf{Rôle} & \textbf{Nom / Entité} & \textbf{Date} & \textbf{Visa} \\
\midrule
Rédacteur & \dotfill & \dotfill & \dotfill \\
Vérificateur & Responsable Infrastructure Réseau & \dotfill & \dotfill \\
Approbateur & Responsable Sécurité & \dotfill & \dotfill \\
\bottomrule
\end{tabular}

\vspace{10pt}
\begin{notebox}[Portée du document]
Ce cahier des charges est le référentiel d'exigences du projet. Les exigences y sont identifiées de manière unique (\rid{BF}, \rid{BNF}, \rid{C}) et rattachées aux cas de recette de la section~8. Le périmètre défini au \S1.1 est figé ; toute évolution fait l'objet d'une nouvelle version.
\end{notebox}

\vspace{6pt}
{\setstretch{1.0}\setcounter{tocdepth}{1}\tableofcontents}

\newpage

%==============================================================================
\section*{Glossaire}
\addcontentsline{toc}{section}{Glossaire}
%==============================================================================
{\small
\begin{longtable}{K L{5.8cm} K L{5.4cm}}
802.1Q & Étiquetage VLAN des trames. & MAB & Authentification par adresse MAC. \\
802.1X & Contrôle d'accès réseau par port. & MAC & Identifiant physique d'interface. \\
ACL & Liste de contrôle d'accès. & NAC & Contrôle d'accès au réseau. \\
CAM & Table des adresses MAC apprises. & NGFW & Pare-feu à inspection applicative. \\
DAI & Contrôle de cohérence ARP. & RADIUS & Service d'authentification réseau. \\
FC & Fibre Channel (réseau de stockage). & SVI & Interface de routage d'un VLAN. \\
IPSG & Filtrage des IP sources illégitimes. & vPC & Agrégation sur deux châssis. \\
LACP & Protocole d'agrégation de liens. & VLAN & Réseau local virtuel. \\
LUN & Volume logique exposé par une baie. & WWPN & Identifiant d'un port Fibre Channel. \\
\end{longtable}}

%==============================================================================
\section{Contexte, problématique et objectifs}
%==============================================================================

\subsection{Objet et périmètre}
Le présent document définit les besoins, l'architecture cible et les critères d'acceptation du projet d'évolution de l'architecture réseau vers un modèle où la politique réseau est attachée à l'\emph{identité} de l'équipement et non à son emplacement physique. Compte tenu de la durée du stage (45 jours), le projet porte sur la \textbf{conception et la validation sur maquette virtualisée}, préparant un déploiement ultérieur.

\begin{table}[H]
\centering
\caption{Périmètre du projet}
\begin{tabularx}{\textwidth}{L{1.7cm} Y}
\toprule
\rowcolor{soft}\textbf{Statut} & \textbf{Éléments} \\
\midrule
\textbf{Inclus} & Conception de l'architecture cible ; maquette virtualisée ; authentification des équipements par 802.1X et MAB avec affectation dynamique de VLAN ; agrégation de liens et répartition de charge par hachage MAC ; durcissement de la couche 2 ; filtrage inter-VLAN sur pare-feu ; raccordement stockage par identité (WWPN) ; supervision ; recette documentée. \\
\midrule
\textbf{Exclu} & Déploiement en production ; réseau WAN et sans fil ; refonte du plan d'adressage de production ; acquisition de matériel ; migration applicative. \\
\bottomrule
\end{tabularx}
\end{table}

\subsection{Architecture existante}
La segmentation logique est aujourd'hui assurée par des VLAN, et l'affectation d'un équipement à son VLAN repose sur une \emph{configuration statique du port physique} sur lequel il est raccordé. Chaque port est associé, dans la documentation d'exploitation, à un identifiant d'équipement. Tout déplacement impose de reconfigurer manuellement le commutateur.

\begin{figure}[H]
\centering
\begin{tikzpicture}[
  font=\footnotesize, scale=0.94, transform shape,
  sw/.style={draw=primary,fill=soft,rounded corners=2pt,minimum width=2.6cm,minimum height=0.75cm,align=center},
  core/.style={draw=primary,fill=primary!85,text=white,rounded corners=2pt,minimum width=3.0cm,minimum height=0.8cm,align=center},
  dev/.style={draw=gris,fill=white,minimum width=1.6cm,minimum height=0.62cm,align=center,font=\scriptsize},
  lnk/.style={-,thick,primary!80},
  blk/.style={-,thick,red!70!black,dashed}]
\node[core] (core) at (0,3.6) {C\oe ur de réseau\\\scriptsize Routage inter-VLAN par ACL};
\node[sw] (a1) at (-3.2,1.9) {Accès SW-A};
\node[sw] (a2) at ( 3.2,1.9) {Accès SW-B};
\draw[lnk] (core) -- (a1); \draw[lnk] (core) -- (a2);
\draw[blk] (a1) -- node[midway,above,font=\scriptsize,text=red!70!black]{lien bloqué par STP} (a2);
\node[dev] (p1) at (-4.4,0.5) {Serveur\\\ttfamily Gi1/0/1};
\node[dev] (p2) at (-2.1,0.5) {Banc test\\\ttfamily Gi1/0/2};
\node[dev] (p3) at ( 2.1,0.5) {Poste\\\ttfamily Gi1/0/5};
\node[dev] (p4) at ( 4.4,0.5) {Imprimante\\\ttfamily Gi1/0/6};
\draw[lnk] (a1) -- (p1); \draw[lnk] (a1) -- (p2);
\draw[lnk] (a2) -- (p3); \draw[lnk] (a2) -- (p4);
\node[draw=warn,fill=warn!8,rounded corners=2pt,align=center,text width=7.4cm,font=\scriptsize]
  at (0,-0.75) {\textbf{Liaison statique} --- 1 port physique $\rightarrow$ 1 VLAN $\rightarrow$ 1 identifiant d'équipement\\
  \textit{Tout déplacement impose une reconfiguration manuelle du commutateur.}};
\end{tikzpicture}
\caption{Architecture existante --- affectation VLAN liée au port physique}
\end{figure}

\subsection{Limites identifiées}
\begin{table}[H]
\centering
\caption{Analyse critique de l'existant}
\begin{tabularx}{\textwidth}{L{1.3cm} L{4.5cm} Y C{1.6cm}}
\toprule
\rowcolor{soft}\textbf{Réf.} & \textbf{Limite} & \textbf{Impact opérationnel} & \textbf{Criticité} \\
\midrule
\rid{L-01} & Couplage port physique / identité de l'équipement & Toute mobilité impose une reconfiguration manuelle (30 à 60 min d'immobilisation). & Élevée \\
\rid{L-02} & Absence d'authentification de l'équipement & Un port actif donne accès au VLAN sans vérifier ce qui s'y raccorde. & Critique \\
\rid{L-03} & Risque d'erreur humaine & Erreur de VLAN ou d'ACL entraînant indisponibilité ou exposition. & Élevée \\
\rid{L-04} & Liens redondants bloqués par STP & La bande passante disponible n'est pas exploitée. & Moyenne \\
\rid{L-05} & Filtrage inter-VLAN limité aux ACL & Ni l'application ni le contenu des flux ne sont inspectés. & Élevée \\
\bottomrule
\end{tabularx}
\end{table}

\begin{notebox}[Problématique]
Comment faire évoluer une architecture dans laquelle l'appartenance d'un équipement à un segment logique est déterminée par son \emph{emplacement physique}, vers une architecture où elle est déterminée par son \emph{identité}, tout en améliorant l'utilisation de la bande passante et le niveau de sécurité ?
\end{notebox}

\subsection{Objectifs}
\begin{table}[H]
\centering
\caption{Objectifs et indicateurs}
\begin{tabularx}{\textwidth}{L{1.4cm} Y L{4.6cm}}
\toprule
\rowcolor{soft}\textbf{Réf.} & \textbf{Objectif} & \textbf{Indicateur / cible} \\
\midrule
\rid{OBJ-01} & Découpler l'identité de l'équipement de son emplacement physique. & 0 action manuelle lors d'un déplacement. \\
\rid{OBJ-02} & Réduire le délai de mise en service et de déplacement. & $\leq$ 30 s au lieu de 30 à 60 min. \\
\rid{OBJ-03} & Authentifier tout équipement avant l'octroi d'un accès. & 100 \% des ports d'accès sous contrôle. \\
\rid{OBJ-04} & Exploiter l'intégralité des liens redondants. & Écart de charge entre liens $\leq$ 20 \%. \\
\rid{OBJ-05} & Inspecter et superviser les flux inter-VLAN. & 100 \% des flux inspectés, détection $\leq$ 60 s. \\
\bottomrule
\end{tabularx}
\end{table}

%==============================================================================
\section{Expression des besoins}
%==============================================================================
Convention : \textbf{DOIT} désigne une exigence impérative conditionnant la recette, \textbf{DEVRAIT} une exigence recommandée dont l'écart doit être justifié.

\subsection{Besoins fonctionnels}
\begin{table}[H]
\centering
\begin{tabularx}{\textwidth}{L{1.4cm} Y C{2.3cm}}
\toprule
\rowcolor{soft}\textbf{Réf.} & \textbf{Besoin fonctionnel} & \textbf{Niveau} \\
\midrule
\rid{BF-01} & Identifier chaque équipement par son adresse MAC lors de son raccordement à un port d'accès. & Impératif \\
\rid{BF-02} & Affecter dynamiquement le VLAN correspondant à l'équipement authentifié, quel que soit le port utilisé. & Impératif \\
\rid{BF-03} & Appliquer dynamiquement la politique de filtrage associée au profil de l'équipement. & Impératif \\
\rid{BF-04} & Authentifier par 802.1X les équipements disposant d'un supplicant, avec repli automatique sur MAB. & Impératif \\
\rid{BF-05} & Placer tout équipement non reconnu dans un VLAN de quarantaine à accès restreint. & Impératif \\
\rid{BF-06} & Répartir le trafic sur tous les liens d'un groupe agrégé selon un hachage des adresses MAC source et destination. & Impératif \\
\rid{BF-07} & Faire transiter l'ensemble des flux inter-VLAN par un pare-feu appliquant des profils de sécurité. & Impératif \\
\rid{BF-08} & Autoriser les accès au stockage sur la base de l'identité de l'initiateur (WWPN) et non du port de raccordement. & Impératif \\
\rid{BF-09} & Superviser l'état des liens, des authentifications et des équipements, et alerter en cas d'anomalie. & Impératif \\
\rid{BF-10} & Produire un inventaire dynamique associant adresse MAC, VLAN attribué, port et horodatage. & Recommandé \\
\bottomrule
\end{tabularx}
\end{table}

\subsection{Besoins non fonctionnels et contraintes}
\begin{table}[H]
\centering
\begin{tabularx}{\textwidth}{L{1.6cm} L{2.7cm} Y}
\toprule
\rowcolor{soft}\textbf{Réf.} & \textbf{Catégorie} & \textbf{Exigence} \\
\midrule
\rid{BNF-01} & Performance & L'authentification et l'affectation du VLAN \textbf{DOIVENT} s'achever en moins de 30 secondes. \\
\rid{BNF-02} & Disponibilité & La perte d'un lien d'un groupe agrégé \textbf{DOIT} être transparente (convergence $\leq$ 1 s). \\
\rid{BNF-03} & Disponibilité & L'indisponibilité du service d'authentification \textbf{NE DOIT PAS} interrompre les sessions établies. \\
\rid{BNF-04} & Sécurité & Aucun accès au réseau \textbf{NE DOIT} être possible sans authentification préalable. \\
\rid{BNF-05} & Sécurité & Les mesures de durcissement de la section~5.2 \textbf{DOIVENT} couvrir 100 \% des ports d'accès. \\
\rid{BNF-06} & Sécurité & Les journaux d'authentification et de filtrage \textbf{DOIVENT} être centralisés et horodatés. \\
\rid{BNF-07} & Exploitabilité & Toute anomalie \textbf{DOIT} être détectée et notifiée en moins de 60 secondes. \\
\rid{BNF-08} & Maintenabilité & Les configurations \textbf{DOIVENT} être versionnées et restaurables. \\
\midrule
\rid{C-01} & Technique & Réutilisation du parc de commutation existant ; aucun renouvellement matériel. \\
\rid{C-02} & Technique & Les équipements de bancs d'essais n'ont pas tous de supplicant 802.1X : MAB est obligatoire pour ceux-ci. \\
\rid{C-03} & Technique & Maquette hébergée sur un poste macOS : ressources limitées, déploiement par lots. \\
\rid{C-04} & Sécurité & Le modèle de sécurité \textbf{NE DOIT PAS} reposer sur la seule adresse MAC (cf. \S5.1). \\
\rid{C-05} & Planning & Réalisation en 45 jours, restitution et rapport inclus. \\
\rid{C-06} & Budgétaire & Priorité aux licences détenues et aux versions d'évaluation ou libres. \\
\bottomrule
\end{tabularx}
\end{table}

%==============================================================================
\section{Architecture cible}
%==============================================================================

\subsection{Principes directeurs}
\begin{enumerate}[leftmargin=*,itemsep=1pt,topsep=2pt]
\item \textbf{L'identité prime sur l'emplacement} : la politique réseau est attachée à l'équipement, pas au port.
\item \textbf{Défense en profondeur} : l'identification par MAC est un mécanisme de \emph{provisionnement}, complété par des contrôles indépendants qui en compensent les faiblesses.
\item \textbf{Point de contrôle unique} : aucun trafic entre VLAN n'échappe à l'inspection.
\item \textbf{Observabilité et réversibilité} : chaque décision est journalisée, chaque lot est retirable indépendamment.
\end{enumerate}

\subsection{Vue d'ensemble}
\begin{figure}[H]
\centering
\begin{tikzpicture}[
  font=\footnotesize, scale=0.93, transform shape,
  box/.style={draw=primary,fill=soft,rounded corners=2pt,minimum width=2.5cm,minimum height=0.75cm,align=center},
  fw/.style={draw=warn,fill=warn!12,rounded corners=2pt,minimum width=2.9cm,minimum height=0.85cm,align=center},
  srv/.style={draw=accent,fill=accent!10,rounded corners=2pt,minimum width=2.4cm,minimum height=0.7cm,align=center,font=\scriptsize},
  dev/.style={draw=gris,fill=white,minimum width=1.6cm,minimum height=0.6cm,align=center,font=\scriptsize},
  lnk/.style={-,thick,primary!80},
  agg/.style={-,very thick,accent}]
\node[fw] (pa) at (0,4.6) {Pare-feu NGFW\\\scriptsize Routage inter-VLAN + profils};
\node[box] (n1) at (-2.5,3.1) {Nexus 1};
\node[box] (n2) at ( 2.5,3.1) {Nexus 2};
\draw[lnk] (pa) -- (n1); \draw[lnk] (pa) -- (n2);
\draw[agg] (n1) -- (n2);
\node[font=\scriptsize,text=accent,fill=white,inner sep=1.5pt] at (0,3.1) {vPC peer-link};
\node[box] (acc) at (0,1.5) {Commutateur d'accès\\\scriptsize 802.1X / MAB};
\draw[agg] (n1) -- node[midway,left,font=\scriptsize,text=accent]{Port-Channel} (acc);
\draw[agg] (n2) -- node[midway,right,font=\scriptsize,text=accent]{src-dst-mac} (acc);
\node[dev] (d1) at (-3.2,0.1) {Serveur};
\node[dev] (d2) at (-1.1,0.1) {Banc test};
\node[dev] (d3) at ( 1.1,0.1) {Poste};
\node[dev] (d4) at ( 3.2,0.1) {Capteur};
\foreach \d in {d1,d2,d3,d4} { \draw[lnk] (acc) -- (\d); }
\node[font=\scriptsize,text=okgreen] at (0,-0.7) {\textbf{Port indifférent} --- la politique suit l'adresse MAC};
\node[srv] (rad) at (-5.6,3.1) {RADIUS\\\scriptsize Référentiel MAC};
\node[srv] (cen) at (-5.6,1.9) {Centreon\\\scriptsize Supervision};
\node[srv] (net) at ( 5.6,3.1) {NetApp\\\scriptsize LUN / zoning WWPN};
\node[srv] (sec) at ( 5.6,1.9) {Scan vuln.\\\scriptsize OpenVAS / ClamAV};
\draw[lnk,dashed] (rad) -- (n1);
\draw[lnk,dashed] (net) -- (n2);
\draw[lnk,dashed] (cen) .. controls (-3.4,1.3) .. (acc);
\draw[lnk,dashed] (sec) .. controls ( 3.4,1.3) .. (acc);
\end{tikzpicture}
\caption{Architecture cible --- provisionnement par identité, agrégation de liens et inspection centralisée}
\end{figure}

\subsection{Identification et provisionnement par adresse MAC}
L'affectation dynamique du VLAN repose sur le couple \textbf{802.1X / MAB} associé à un serveur \textbf{RADIUS}. Le commutateur joue le rôle d'authentificateur : tant qu'aucune authentification n'a abouti, le port ne laisse passer que le trafic d'authentification.

\begin{itemize}[leftmargin=*,itemsep=1pt,topsep=2pt]
\item \textbf{802.1X} en priorité pour les équipements disposant d'un supplicant (postes, serveurs) : l'identité est portée par un certificat, ce qui offre la meilleure garantie.
\item \textbf{MAB} pour les équipements sans supplicant (bancs d'essais, capteurs, imprimantes) : le commutateur apprend l'adresse MAC de la première trame et la soumet au RADIUS.
\item \textbf{Quarantaine} (VLAN 99) en cas d'échec, avec pour seuls flux autorisés DHCP, DNS et un portail d'information.
\end{itemize}

\begin{notebox}[Clarification technique]
L'algorithme \ttb{port-channel load-balance src-dst-mac} réalise une \emph{répartition de charge} entre les liens d'un groupe agrégé. Il ne supprime pas, à lui seul, la reconfiguration lors d'un déplacement d'équipement : cette propriété est apportée par le couple \textbf{MAB/802.1X + affectation dynamique de VLAN par RADIUS}. Les deux mécanismes sont complémentaires et tous deux retenus : le premier pour la performance (\rid{OBJ-04}), le second pour la mobilité (\rid{OBJ-01}, \rid{OBJ-02}).
\end{notebox}

\begin{figure}[H]
\centering
\begin{tikzpicture}[font=\scriptsize, scale=0.92, transform shape,
  lane/.style={draw=primary,fill=soft,minimum width=2.4cm,minimum height=0.6cm,align=center,font=\footnotesize\bfseries},
  msg/.style={-{Latex[length=2mm]},thick,primary!85},
  act/.style={draw=primary,fill=soft,rounded corners=2pt,font=\scriptsize,inner sep=3pt},
  lbl/.style={above,font=\scriptsize,midway}]
\node[lane] (E) at (0,0)    {Équipement};
\node[lane] (S) at (4.4,0)  {Commutateur};
\node[lane] (R) at (8.8,0)  {RADIUS};
\node[lane] (N) at (12.8,0) {Réseau / VLAN};
\foreach \n in {E,S,R,N} { \draw[dashed,gris] (\n.south) -- ++(0,-6.1); }
\draw[msg] (0,-0.8) -- (4.4,-0.8) node[lbl]{1. Raccordement, première trame};
\draw[msg] (4.4,-1.6) -- (8.8,-1.6) node[lbl]{2. Access-Request (MAC / EAP)};
\node[act] at (8.8,-2.3) {3. Vérification dans le référentiel};
\draw[msg] (8.8,-3.1) -- (4.4,-3.1) node[lbl]{4. Access-Accept + VLAN + Filter-Id};
\node[act] at (4.4,-3.8) {5. Application du VLAN et de l'ACL sur le port};
\draw[msg] (4.4,-4.6) -- (12.8,-4.6) node[lbl]{6. Accès autorisé, session journalisée};
\draw[msg,red!70!black,dashed] (4.4,-5.7) -- (12.8,-5.7)
  node[lbl,text=red!70!black]{4bis. En cas de rejet : bascule en VLAN 99 (quarantaine)};
\end{tikzpicture}
\caption{Séquence d'authentification MAB et d'affectation dynamique de VLAN}
\end{figure}

Les attributs renvoyés par le serveur sont conformes à la RFC 3580 : \ttb{Tunnel-Type} (valeur VLAN), \ttb{Tunnel-Medium-Type} (valeur IEEE-802), \ttb{Tunnel-Private-Group-Id} portant l'identifiant du VLAN à appliquer, et \ttb{Filter-Id} portant le nom de l'ACL associée au profil. En cas d'indisponibilité du RADIUS, les sessions établies sont maintenues et un VLAN de secours est appliqué aux nouveaux raccordements (\rid{BNF-03}).

\subsection{Répartition de charge par hachage MAC}
Les liens entre l'accès et le c\oe ur sont agrégés en \textbf{Port-Channel} négocié par \textbf{LACP}. L'algorithme \ttb{src-dst-mac} calcule une empreinte à partir des adresses MAC source et destination ; le résultat détermine le lien emprunté. Il en découle un \emph{déterminisme par flux} (aucun déséquencement de paquets) et une \emph{répartition statistique} d'autant meilleure que les couples d'adresses sont variés --- ce qui correspond au profil de trafic de la plateforme. Sur les commutateurs Nexus, la fonction \textbf{vPC} répartit les liens d'un même Port-Channel sur deux châssis : la redondance ne repose plus sur le blocage de liens par STP (\rid{L-04}).

\begin{alertbox}[Limite du hachage par MAC]
La granularité de répartition est le \emph{flux}, jamais le paquet : un unique flux entre deux équipements ne dépassera pas le débit d'un lien physique. Si le trafic était dominé par quelques flux à très haut débit, l'algorithme \ttb{src-dst-ip} devrait être évalué. Ce point fait l'objet du cas de test \rid{TC-06}.
\end{alertbox}

\subsection{Segmentation, routage inter-VLAN et stockage}
La segmentation reste assurée par les VLAN 802.1Q. La différence porte sur le \emph{point de routage} : le routage inter-VLAN, aujourd'hui réalisé par des SVI filtrées par ACL, est transféré au pare-feu NGFW raccordé en trunk avec une sous-interface de niveau 3 par VLAN. Chaque VLAN devient une zone de sécurité, et tout flux entre zones traverse le moteur de politique (\rid{OBJ-05}).

\begin{notebox}[Fibre Channel et adresses MAC]
Le protocole Fibre Channel ne transporte pas d'adresses MAC : il identifie les ports par un \textbf{WWPN} sur 64 bits. Le principe est néanmoins identique et confirme la démarche : le \emph{zoning logiciel} (par WWPN) attache l'autorisation à l'identité de l'initiateur, tandis que le \emph{zoning matériel} (par port de fabrique) l'attache à l'emplacement. Déplacer un serveur d'un port à un autre est transparent dans le premier cas, et impose une reconfiguration dans le second --- exactement la problématique traitée côté LAN.
\end{notebox}

Côté NetApp, le zoning \textbf{DOIT} donc être réalisé par WWPN (un initiateur, une cible par zone), le masquage des LUN s'appuyer sur des \textit{igroups} construits sur ces WWPN, et les flux de stockage IP être confinés dans un VLAN dédié non routé vers les VLAN utilisateurs.

\subsection{Supervision et comparatif}
Centreon assure la collecte et la notification : état de chaque membre du Port-Channel et écart de charge, disponibilité du RADIUS et taux d'échec d'authentification, adresses en quarantaine, violations de \textit{port-security} et MAC \textit{flapping}, charge du pare-feu, état des chemins de stockage et des VM.

\begin{table}[H]
\centering
\caption{Synthèse comparative}
\begin{tabularx}{\textwidth}{L{3.5cm} Y Y}
\toprule
\rowcolor{soft}\textbf{Critère} & \textbf{Existant} & \textbf{Cible} \\
\midrule
Affectation du VLAN & Statique, liée au port & Dynamique, liée à l'adresse MAC \\
Déplacement d'un équipement & Reconfiguration manuelle & Automatique, sans intervention \\
Authentification & Aucune & 802.1X, repli MAB, quarantaine \\
Redondance des liens & STP, lien bloqué & LACP / vPC, tous liens actifs \\
Filtrage inter-VLAN & ACL (L3/L4) & NGFW à inspection applicative \\
Accès stockage & Selon le port de fabrique & Selon le WWPN de l'initiateur \\
Inventaire et traçabilité & Documentaire, partielle & Dynamique, journalisée \\
\bottomrule
\end{tabularx}
\end{table}

%==============================================================================
\section{Exigences de sécurité}
%==============================================================================

\subsection{Risques et parades}
\begin{table}[H]
\centering
\begin{tabularx}{\textwidth}{L{1.3cm} L{3.9cm} C{1.5cm} Y}
\toprule
\rowcolor{soft}\textbf{Réf.} & \textbf{Risque} & \textbf{Criticité} & \textbf{Mesure de réduction} \\
\midrule
\rid{R-01} & Usurpation d'adresse MAC & Critique & 802.1X prioritaire ; \textit{port-security} ; DHCP Snooping + DAI + IPSG ; alerte sur MAC dupliquée ou \textit{flapping}. \\
\rid{R-02} & Indisponibilité du RADIUS & Élevée & Deux instances ; maintien des sessions établies ; VLAN de secours. \\
\rid{R-03} & VLAN hopping (double étiquetage, DTP) & Élevée & VLAN natif dédié et inutilisé ; DTP désactivé ; \textit{pruning} des VLAN sur les trunks. \\
\rid{R-04} & Saturation de la table CAM & Élevée & \textit{port-security maximum} ; \textit{storm-control} ; supervision du remplissage. \\
\rid{R-05} & Boucle de niveau 2 & Moyenne & BPDU Guard, Root Guard, PortFast sur les ports d'accès. \\
\rid{R-06} & Vulnérabilité non corrigée & Élevée & Scan hebdomadaire, plan de correction, durcissement des images. \\
\bottomrule
\end{tabularx}
\end{table}

\begin{alertbox}[Le point de sécurité central du projet]
Une adresse MAC est lisible sur le réseau et modifiable logiciellement : c'est un \emph{identifiant}, jamais un \emph{secret}. Le MAB ne constitue donc pas une authentification forte. Il est retenu ici comme mécanisme de \textbf{provisionnement automatisé}, obligatoirement adossé à trois contrôles indépendants : 802.1X partout où il est possible (\rid{BF-04}), le durcissement de la couche 2 (\S5.2), et l'inspection applicative de tous les flux inter-VLAN (\rid{BF-07}). Cette exigence est formalisée par la contrainte \rid{C-04}.
\end{alertbox}

\subsection{Durcissement de la couche 2}
\begin{table}[H]
\centering
\begin{tabularx}{\textwidth}{L{4.6cm} Y C{2.0cm}}
\toprule
\rowcolor{soft}\textbf{Mesure} & \textbf{Objectif} & \textbf{Risque couvert} \\
\midrule
\textit{Port-security} (1 à 2 MAC, mode \textit{restrict}) & Limiter les adresses apprises par port et alerter en cas de dépassement. & R-01, R-04 \\
DHCP Snooping + DAI + IP Source Guard & Construire la table IP/MAC de confiance et rejeter ARP et IP incohérents. & R-01 \\
VLAN natif dédié, DTP désactivé, \textit{pruning} & Neutraliser le double étiquetage et la négociation de trunk. & R-03 \\
BPDU Guard / Root Guard / PortFast & Protéger la topologie spanning-tree. & R-05 \\
\textit{Storm-control} & Limiter les tempêtes de diffusion. & R-04 \\
Désactivation des ports inutilisés & Réduire la surface d'exposition physique. & R-01 \\
\bottomrule
\end{tabularx}
\end{table}

\subsection{Filtrage des flux et inspection}
Le principe appliqué est le \emph{refus par défaut} : tout flux non explicitement autorisé est bloqué et journalisé.

\begin{table}[H]
\centering
\caption{Matrice de flux inter-VLAN (référence)}
\begin{tabularx}{\textwidth}{L{2.6cm} L{2.1cm} L{3.4cm} C{1.7cm} Y}
\toprule
\rowcolor{soft}\textbf{Source} & \textbf{Destination} & \textbf{Service} & \textbf{Action} & \textbf{Profil} \\
\midrule
MGMT & Tous & SSH/22, HTTPS/443, SNMP/161 & Autoriser & Strict \\
USERS & SERVERS & HTTPS/443, HTTP/8080 & Autoriser & Complet \\
BANCS & SERVERS & TCP/5000--5010 & Autoriser & Complet \\
SERVERS & STORAGE & NFS/2049, iSCSI/3260 & Autoriser & Antivirus \\
QUARANTINE & Portail & DHCP/67, DNS/53, HTTP/80 & Autoriser & Journalisé \\
Tous & MGMT & Tout & \textbf{Refuser} & Journalisé \\
Tous & Tous & Tout autre & \textbf{Refuser} & Journalisé \\
\bottomrule
\end{tabularx}
\end{table}

Les règles d'autorisation \textbf{DOIVENT} être associées à un groupe de profils appliquant : antivirus (réinitialisation), anti-spyware avec \textit{sinkhole} DNS, protection contre les vulnérabilités (blocage des sévérités critique et élevée), filtrage d'URL, blocage des types de fichiers exécutables non nécessaires et analyse comportementale des fichiers inconnus. En complément, un scan de vulnérabilités authentifié (OpenVAS) est exécuté chaque semaine avec pour exigence \emph{aucune vulnérabilité critique non traitée}, un scan de découverte réseau détecte tout équipement non inventorié, et une analyse antivirus quotidienne (ClamAV) couvre les serveurs. L'ensemble des équipements exporte ses journaux vers un collecteur syslog centralisé et synchronisé par NTP (\rid{BNF-06}).

%==============================================================================
\section{Environnement technique}
%==============================================================================

\subsection{Plateforme et choix technologiques}
La validation est réalisée sur une maquette virtualisée hébergée sur un poste macOS. L'hyperviseur doit supporter les réseaux virtuels multiples avec étiquetage 802.1Q et le mode promiscuité.

\begin{table}[H]
\centering
\begin{tabularx}{\textwidth}{L{3.0cm} L{3.6cm} Y}
\toprule
\rowcolor{soft}\textbf{Composant} & \textbf{Solution retenue} & \textbf{Justification} \\
\midrule
Hyperviseur macOS & VMware Fusion & Support 802.1Q et réseaux personnalisés ; alternative : UTM/QEMU. \\
Émulation réseau & GNS3 ou EVE-NG & Images Cisco IOSvL2 / Nexus 9000v pour 802.1X, MAB et vPC. \\
Commutation & Cisco Catalyst / Nexus & Parc existant (\rid{C-01}) ; support natif de MAB, vPC et \ttb{src-dst-mac}. \\
Pare-feu & Palo Alto (NGFW) & Identification applicative et profils de sécurité intégrés ; FortiGate en secours. \\
Authentification & FreeRADIUS & Libre, conforme RFC 2865/3580 ; migration ultérieure vers un NAC commercial possible. \\
Stockage & NetApp ONTAP & Zoning WWPN, igroups et multipathing ; conforme à \rid{BF-08}. \\
Supervision & Centreon & Solution déjà exploitée, connecteurs SNMP matures. \\
Systèmes & Ubuntu Server 24.04 LTS & Support long terme, image de référence durcissable. \\
\bottomrule
\end{tabularx}
\end{table}

\begin{alertbox}[Contrainte de ressources]
Le cumul des machines virtuelles dépasse les capacités d'un poste standard. Le déploiement \textbf{DOIT} être réalisé par lots, avec au maximum quatre à cinq machines actives simultanément. Un hôte disposant d'au moins 32~Go de mémoire vive est requis (\rid{C-03}).
\end{alertbox}

\subsection{Machines virtuelles et plan VLAN}
\begin{table}[H]
\centering
\begin{tabularx}{\textwidth}{L{1.4cm} L{2.4cm} Y C{1.0cm} C{1.0cm} C{1.25cm}}
\toprule
\rowcolor{soft}\textbf{Réf.} & \textbf{Nom} & \textbf{Rôle} & \textbf{vCPU} & \textbf{RAM} & \textbf{Disque} \\
\midrule
VM-01 & fw-ngfw-01 & Pare-feu, routage inter-VLAN, profils de sécurité & 4 & 8 Go & 60 Go \\
VM-02 & rad-01 & Ubuntu Server + FreeRADIUS, référentiel MAC & 2 & 4 Go & 25 Go \\
VM-03 & sup-01 & Ubuntu Server + Centreon, OpenVAS, ClamAV & 4 & 8 Go & 70 Go \\
VM-04 & srv-01 & Ubuntu Server applicatif, client stockage & 2 & 4 Go & 40 Go \\
VM-05 & stor-01 & Simulateur NetApp ONTAP, LUN et NFS & 4 & 8 Go & 100 Go \\
VM-06 & clt-01 & Client Ubuntu, supplicant 802.1X & 2 & 4 Go & 25 Go \\
SW-01/03 & sw-nx-01/02, sw-acc-01 & Commutateurs virtuels (vPC, Port-Channel, MAB) & 2 & 4 Go & 8 Go \\
\bottomrule
\end{tabularx}
\end{table}

\begin{table}[H]
\centering
\begin{tabularx}{\textwidth}{C{1.5cm} L{2.6cm} L{3.2cm} L{2.6cm} Y}
\toprule
\rowcolor{soft}\textbf{VLAN} & \textbf{Nom} & \textbf{Sous-réseau} & \textbf{Passerelle} & \textbf{Usage} \\
\midrule
10 & MGMT & 10.10.10.0/24 & 10.10.10.1 & Administration, supervision \\
20 & SERVERS & 10.10.20.0/24 & 10.10.20.1 & Serveurs applicatifs \\
30 & BANCS & 10.10.30.0/24 & 10.10.30.1 & Bancs d'essais, équipements mobiles \\
40 & STORAGE & 10.10.40.0/24 & 10.10.40.1 & Flux de stockage IP \\
50 & USERS & 10.10.50.0/24 & 10.10.50.1 & Postes de travail \\
99 & QUARANTINE & 10.10.99.0/24 & 10.10.99.1 & Équipements non authentifiés \\
999 & NATIVE-UNUSED & --- & --- & VLAN natif des trunks, non utilisé \\
\bottomrule
\end{tabularx}
\end{table}

%==============================================================================
\section{Lotissement et planning}
%==============================================================================
Le projet est découpé en six lots réalisables et réversibles indépendamment, répartis sur les 45 jours du stage.

\begin{table}[H]
\centering
\caption{Lots et jalons}
\begin{tabularx}{\textwidth}{L{1.2cm} L{3.3cm} Y C{1.5cm} C{1.2cm}}
\toprule
\rowcolor{soft}\textbf{Lot} & \textbf{Intitulé} & \textbf{Contenu} & \textbf{Semaine} & \textbf{Jalon} \\
\midrule
L0 & Cadrage & Étude de l'existant, relevé des configurations, validation du cahier des charges. & S1 & J1 \\
L1 & Socle de maquette & Hyperviseur, VM, commutateurs virtuels, plan VLAN et adressage. & S2 & J2 \\
L2 & Authentification & FreeRADIUS, référentiel MAC, 802.1X et MAB, VLAN dynamique, quarantaine. & S3--S4 & J3 \\
L3 & Agrégation & LACP, Port-Channel, hachage \ttb{src-dst-mac}, vPC, mesures de charge. & S4--S5 & --- \\
L4 & Sécurité & Durcissement L2, pare-feu et profils, matrice de flux, scans. & S5--S6 & J4 \\
L5 & Stockage, supervision et recette & Zoning WWPN, sondes Centreon, exécution du cahier de recette, rapport. & S6--S7 & J5 \\
\bottomrule
\end{tabularx}
\end{table}

\begin{figure}[H]
\centering
\begin{ganttchart}[
  hgrid, vgrid,
  x unit=1.05cm, y unit title=0.6cm, y unit chart=0.54cm,
  title height=1, title label font=\scriptsize,
  bar height=0.5, bar label font=\scriptsize,
  group label font=\bfseries\scriptsize,
  milestone label font=\scriptsize\itshape,
  bar/.style={fill=accent!55, draw=accent!85},
  group/.style={fill=primary!85, draw=primary},
  milestone/.style={fill=red!70!black, draw=red!70!black}
]{1}{7}
\gantttitle{Semaines --- 45 jours}{7} \\
\gantttitlelist{1,...,7}{1} \\
\ganttgroup{L0 --- Cadrage}{1}{1} \\
\ganttbar{Étude de l'existant}{1}{1} \\
\ganttbar{Cahier des charges}{1}{1} \\
\ganttmilestone{J1}{1} \\
\ganttgroup{L1 --- Socle de maquette}{2}{2} \\
\ganttbar{Hyperviseur, VM, VLAN}{2}{2} \\
\ganttmilestone{J2}{2} \\
\ganttgroup{L2 --- Authentification}{3}{4} \\
\ganttbar{RADIUS et référentiel MAC}{3}{3} \\
\ganttbar{802.1X / MAB / VLAN dynamique}{3}{4} \\
\ganttmilestone{J3}{4} \\
\ganttgroup{L3 --- Agrégation}{4}{5} \\
\ganttbar{LACP, vPC, hachage src-dst-mac}{4}{5} \\
\ganttgroup{L4 --- Sécurité}{5}{6} \\
\ganttbar{Durcissement L2 et pare-feu}{5}{6} \\
\ganttmilestone{J4}{6} \\
\ganttgroup{L5 --- Stockage, supervision, recette}{6}{7} \\
\ganttbar{NetApp (WWPN) et Centreon}{6}{6} \\
\ganttbar{Recette et rapport de stage}{6}{7} \\
\ganttmilestone{J5}{7}
\end{ganttchart}
\caption{Planning prévisionnel sur 7 semaines (45 jours)}
\end{figure}

\begin{table}[H]
\centering
\caption{Livrables}
\begin{tabularx}{\textwidth}{L{1.4cm} L{4.6cm} Y C{1.5cm}}
\toprule
\rowcolor{soft}\textbf{Réf.} & \textbf{Livrable} & \textbf{Description} & \textbf{Jalon} \\
\midrule
LV-01 & Cahier des charges & Le présent document validé. & J1 \\
LV-02 & Dossier d'architecture & Schémas, plan VLAN, matrice de flux. & J2 \\
LV-03 & Configurations de référence & Fichiers versionnés, commentés et restaurables. & J4 \\
LV-04 & Rapport de tests & Mesures de délai, de répartition de charge et de sécurité. & J4 \\
LV-05 & PV de recette et rapport de stage & Résultats TC-01 à TC-10, procédures d'exploitation, soutenance. & J5 \\
\bottomrule
\end{tabularx}
\end{table}

Les principaux risques projet sont les ressources insuffisantes sur l'hôte de maquettage (parade : déploiement par lots, arbitrage sur un serveur mutualisé), l'indisponibilité d'images ou de licences (parade : anticipation dès S1, substituts libres identifiés) et les fonctionnalités non supportées par les commutateurs virtuels (parade : validation dès L1).

%==============================================================================
\section{Tests et recette}
%==============================================================================

\begin{table}[H]
\centering
\caption{Cahier de recette}
\begin{tabularx}{\textwidth}{L{1.4cm} L{4.6cm} Y}
\toprule
\rowcolor{soft}\textbf{Réf.} & \textbf{Objet du test} & \textbf{Critère d'acceptation} \\
\midrule
\rid{TC-01} & Authentification 802.1X d'un poste & Accès accordé, VLAN attendu affecté, session journalisée. \\
\rid{TC-02} & Authentification MAB d'un équipement sans supplicant & Accès accordé sur la base de la MAC en moins de 30 s. \\
\rid{TC-03} & Déplacement vers un autre port, puis un autre commutateur & VLAN et ACL identiques, aucune action manuelle, délai $\leq$ 30 s. \\
\rid{TC-04} & Raccordement d'un équipement inconnu & Placement en VLAN 99, accès restreint, alerte remontée. \\
\rid{TC-05} & Indisponibilité du serveur RADIUS & Sessions maintenues, nouveaux accès en VLAN de secours, alerte émise. \\
\rid{TC-06} & Répartition de charge \ttb{src-dst-mac} & Trafic réparti sur tous les membres, écart $\leq$ 20 \% avec $\geq$ 8 couples de MAC. \\
\rid{TC-07} & Perte d'un membre du Port-Channel ou d'un châssis vPC & Continuité du trafic, convergence $\leq$ 1 s, aucune session TCP perdue. \\
\rid{TC-08} & Filtrage inter-VLAN et profils de sécurité & Flux autorisés fonctionnels, flux non autorisés bloqués et journalisés ; fichier de test EICAR bloqué. \\
\rid{TC-09} & Usurpation d'adresse MAC et saturation CAM & Détection par \textit{port-security} / DAI, blocage ou alerte, journalisation. \\
\rid{TC-10} & Accès stockage après changement de port de fabrique, et supervision & LUN toujours accessible sans modifier le zoning ; toute anomalie notifiée en $\leq$ 60 s. \\
\bottomrule
\end{tabularx}
\end{table}

\begin{table}[H]
\centering
\caption{Critères d'acceptation globaux}
\begin{tabularx}{\textwidth}{Y C{3.0cm} C{3.0cm}}
\toprule
\rowcolor{soft}\textbf{Indicateur} & \textbf{Cible} & \textbf{Seuil minimal} \\
\midrule
Délai de mise en service après déplacement & $\leq$ 30 s & $\leq$ 60 s \\
Actions manuelles par déplacement & 0 & 0 \\
Taux d'authentification aboutie & $\geq$ 99 \% & $\geq$ 97 \% \\
Écart de charge entre membres du Port-Channel & $\leq$ 10 \% & $\leq$ 20 \% \\
Convergence après perte de lien & $\leq$ 1 s & $\leq$ 3 s \\
Couverture d'inspection des flux inter-VLAN & 100 \% & 100 \% \\
Vulnérabilités critiques non traitées & 0 & 0 \\
Cas de recette impératifs réussis & 100 \% & 100 \% \\
\bottomrule
\end{tabularx}
\end{table}

%==============================================================================
\section{Conclusion}
%==============================================================================
Le projet substitue, dans la chaîne de décision réseau, l'\emph{identité de l'équipement} à son \emph{emplacement physique}. Cette substitution repose sur trois piliers complémentaires : l'affectation dynamique du VLAN par authentification 802.1X ou MAB via RADIUS, qui supprime la reconfiguration liée à la mobilité ; l'agrégation de liens avec répartition de charge par hachage des adresses MAC, qui valorise toute la bande passante disponible ; et le transfert du routage inter-VLAN vers un pare-feu à inspection applicative. Le même principe est décliné côté stockage par un zoning fondé sur les WWPN plutôt que sur les ports de fabrique.

Le gain attendu est mesurable : suppression des actions manuelles lors d'un déplacement, réduction du délai de mise en service de plusieurs dizaines de minutes à moins d'une minute, et exploitation complète des liens redondants.

L'atout de cette architecture est aussi son principal point de vigilance : une adresse MAC est un identifiant public et falsifiable. C'est pourquoi le MAB n'est jamais employé seul, mais toujours adossé à 802.1X lorsque c'est possible, au durcissement de la couche 2 et à un filtrage applicatif en profondeur. Sous cette réserve, formalisée par la contrainte \rid{C-04}, l'architecture cible répond aux objectifs \rid{OBJ-01} à \rid{OBJ-05} et constitue un socle validé pour un déploiement ultérieur en production.

%==============================================================================
\appendix
\section{Annexe --- Configurations de référence}
%==============================================================================

\subsection{Commutateur d'accès --- 802.1X, MAB et durcissement}
\begin{lstlisting}
! --- Parametres AAA / RADIUS ---
aaa new-model
aaa authentication dot1x default group radius
aaa authorization network default group radius
dot1x system-auth-control
radius server RAD-01
 address ipv4 10.10.10.20 auth-port 1812 acct-port 1813
 key <CLE_PARTAGEE>
!
! --- Port d'acces : 802.1X prioritaire, repli MAB, quarantaine ---
interface range GigabitEthernet1/0/1 - 24
 switchport mode access
 switchport nonegotiate
 authentication order dot1x mab
 authentication priority dot1x mab
 authentication port-control auto
 authentication event fail action authorize vlan 99
 authentication event server dead action authorize vlan 99
 mab
 dot1x pae authenticator
 switchport port-security
 switchport port-security maximum 2
 switchport port-security violation restrict
 spanning-tree portfast
 spanning-tree bpduguard enable
!
! --- Durcissement de couche 2 ---
ip dhcp snooping
ip dhcp snooping vlan 10,20,30,40,50
ip arp inspection vlan 10,20,30,40,50
ip verify source
!
! --- Verification ---
show authentication sessions interface Gi1/0/3 details
\end{lstlisting}

\subsection{Agrégation de liens et répartition de charge}
\begin{lstlisting}
! --- Algorithme de hachage global base sur les adresses MAC ---
port-channel load-balance src-dst-mac
!
interface Port-channel1
 switchport mode trunk
 switchport trunk native vlan 999
 switchport trunk allowed vlan 10,20,30,40,50,99
!
interface range TenGigabitEthernet1/1/1 - 2
 channel-protocol lacp
 channel-group 1 mode active
!
! --- Verification ---
show etherchannel load-balance
show etherchannel 1 summary
\end{lstlisting}

\subsection{FreeRADIUS --- référentiel des équipements (MAB)}
\begin{lstlisting}
# /etc/freeradius/3.0/mods-config/files/authorize
# Adresse MAC en minuscules, sans separateur

0050569a1c3f  Cleartext-Password := "0050569a1c3f"
    Tunnel-Type = VLAN,
    Tunnel-Medium-Type = IEEE-802,
    Tunnel-Private-Group-Id = "30",
    Filter-Id = "ACL-BANCS"

# Refus par defaut : tout equipement inconnu part en quarantaine
DEFAULT  Auth-Type := Reject
\end{lstlisting}

\subsection{Stockage --- autorisation par identité (WWPN)}
\begin{lstlisting}
# Groupe d'initiateurs base sur le WWPN, independant du port de fabrique
igroup create -vserver svm_lab -igroup ig_srv_01 -protocol fcp -ostype linux -initiator 21:00:00:24:ff:8b:1a:2c
lun map -vserver svm_lab -path /vol/vol_app/lun_app01 -igroup ig_srv_01 -lun-id 0

# Zoning logiciel de la fabrique : un initiateur, une cible
zonecreate "z_srv01_netapp", "21:00:00:24:ff:8b:1a:2c; 20:01:00:a0:98:11:22:33"
cfgadd "cfg_prod", "z_srv01_netapp"
cfgenable "cfg_prod"
\end{lstlisting}

\end{document}
